% Generated draft: requires mathematical and source review.
\documentclass[11pt]{article}
\usepackage[margin=25mm]{geometry}
\usepackage{fontspec}
\setmainfont{Latin Modern Roman}
\usepackage{amsmath,amssymb,mathtools,unicode-math}
\setmathfont{Latin Modern Math}
\usepackage{enumitem,needspace,xurl,hyperref}
\setlength{\parindent}{0pt}
\setlength{\parskip}{5pt}
\setlength{\emergencystretch}{3em}
\newcommand{\N}{\mathbb{N}}
\newcommand{\Z}{\mathbb{Z}}
\newcommand{\Q}{\mathbb{Q}}
\newcommand{\R}{\mathbb{R}}
\newcommand{\C}{\mathbb{C}}
\newcommand{\F}{\mathbb{F}}
\newcommand{\A}{\mathbb{A}}
\newcommand{\PP}{\mathbb{P}}
\newcommand{\E}{\mathbb{E}}
\newcommand{\eps}{\varepsilon}
\newcommand{\dd}{\,\mathrm{d}}
\begin{document}
\section*{UnsolvedMath ID: 701061}
\textbf{Source problem code:} catalogue entry
\subsection{Definitions and mathematical statement}
For a prime \(p\), a finite \(p\)-group is a finite group of order \(p^n\) for an integer \(n\geq 0\). A section of a group \(P\) is a quotient \(B/N\), where \(B\leq P\) and \(N\trianglelefteq B\). It is abelian precisely when \([B,B]\leq N\).

The question attributed to L. Pyber asks whether there is an absolute constant \(c>0\), independent of both \(p\) and \(P\), such that
\[
\forall p\text{ prime}\;\forall P\text{ a finite }p\text{-group}\;\exists B\leq P\;\exists N\trianglelefteq B:
\qquad [B,B]\leq N\quad\text{and}\quad |B/N|^c>|P|.
\]
Here positivity of \(c\) expresses the intended size-bound interpretation of the source's phrase ``absolute constant.''

\textit{Unresolved boundary qualification.} The supplied statement does not exclude the trivial group. Under the definition above it includes \(|P|=1\), whose only section has order \(1\), so the strict displayed inequality cannot hold in that case. It requires clarification whether the intended quantifier is over nontrivial finite \(p\)-groups. That restriction has not been silently inserted.

The source reports that a result of A. Yu. Olshanskii prevents requiring the abelian section \(A=B/N\) to be a subgroup of \(P\). It also reports that a result of J. G. Thompson would turn the asserted existence of such sections into the existence of a ``class \(2\)'' subgroup \(H\leq P\) satisfying \(|H|^c>|P|\). These are supplied background qualifications, not independently verified assertions here. The source does not specify whether ``class \(2\)'' means nilpotency class exactly \(2\) or at most \(2\), or explicitly clarify the relationship between the constants in the section and subgroup bounds.
\subsection{Short English statement}
L. Pyber asks whether every finite \(p\)-group has an abelian quotient of some subgroup whose order, raised to one fixed positive power, exceeds the order of the whole group. The same power must work for every prime and every group. The trivial-group boundary needs clarification. The supplied registry labels the problem open; this is not an independently verified status assessment.
\subsection{Sources}
\begin{itemize}
\item Local UnsolvedMath registry, permanent ID \(701061\), exact source code \texttt{catalogue entry}, title \emph{Kourovka Notebook Problem 14.76}. Supplied upstream revision: \texttt{d371ee917ac9647fbb06cd6fc529deec6f3dda82}. The record was created and updated on 2026-10-16 and supplies the status ``open.''
\item \emph{The Kourovka Notebook (Unsolved Problems in Group Theory)}, Problem 14.76, proposed by L. Pyber in the 14th issue (1999); the supplied record identifies the 21st issue edition (2026), \texttt{arXiv:1401.0300v48}. The collection description names E. I. Khukhro and V. D. Mazurov as editors. These bibliographic details are supplied attribution and have not been independently verified.
\item A. Yu. Olshanskii, \emph{Math. Notes} \textbf{23} (1978), 183--185. Cited by the supplied statement for the qualification that one cannot require the abelian section to be a subgroup. No article title is supplied.
\item J. G. Thompson, \emph{J. Algebra} \textbf{13} (1969), 149--151. Cited by the supplied statement for the conditional implication concerning a class \(2\) subgroup with a corresponding power-size bound. No article title is supplied.
\end{itemize}
\end{document}
